\section{Teorema de corte de Fourier}

El teorema de corte de Fourier conecta la transformada de Radon con la transformada de Fourier bidimensional. Con las convenciones del capítulo anterior, el resultado toma una forma directa.

\begin{theorem}[Teorema de corte de Fourier]
Sea $f\in\mathcal S(\mathbb R^2)$. Entonces, para todo $\theta\in[0,\pi)$ y $\sigma\in\mathbb R$,
\[
\mathcal F_t(\mathcal Rf)(\sigma,\theta)
=
\widehat f(\sigma\omega(\theta)).
\]
\end{theorem}

\begin{proof}
Usando la definición de $\mathcal Rf$ y el cambio de variables $x=t\omega(\theta)+s\omega^\perp(\theta)$, se tiene
\[
\begin{aligned}
\mathcal F_t(\mathcal Rf)(\sigma,\theta)
&=
\int_{\mathbb R}
\left(
\int_{\mathbb R}
f(t\omega(\theta)+s\omega^\perp(\theta))\,ds
\right)
e^{-2\pi i t\sigma}\,dt \\
&=
\int_{\mathbb R^2}
f(x)e^{-2\pi i\sigma x\cdot\omega(\theta)}\,dx \\
&=
\widehat f(\sigma\omega(\theta)).
\end{aligned}
\]
\end{proof}

Este resultado muestra que la transformada de Fourier unidimensional de una proyección proporciona los valores de $\widehat f$ sobre la recta radial
\[
\{\sigma\omega(\theta):\sigma\in\mathbb R\}
\]
en el dominio de frecuencias. Al variar $\theta\in[0,\pi)$, estas rectas cubren todo $\mathbb R^2$, aunque la parametrización es degenerada en el origen. El teorema de corte se encuentra, con posibles diferencias de normalización, en Helgason \cite{helgason1999}, Kuchment \cite[Sec.~5.5.4, Thm.~5.10, pp.~32--33]{kuchment2014} y Natterer \cite[Sec.~II.1, Thm.~1.1, p.~11]{natterer2001}; Kak y Slaney lo presentan en el contexto de proyecciones paralelas \cite[Sec.~3.2]{kak_slaney1988}.
